\documentclass[10pt,twocolumn]{article}
\usepackage[margin=0.8in]{geometry}
\usepackage{amsmath,amssymb,booktabs,xcolor}
\usepackage[hidelinks]{hyperref}
\usepackage{caption}
\captionsetup{font=small,labelfont=bf}
% Generated by generate_tables.py -- do not edit by hand.
\newcommand{\ArmRows}{ibm\_marrakesh & B1 probe best & 0.3133 (0.0088) & 0.574 & q116: 3.6 / 4.7\\
 & B2 calibration best & 0.0631 (0.0047) & 0.676 & q46: 34.4 / 1.1\\
 & B3 stock level 3 & 0.0736 (0.0034) & 0.729 & q46: 34.4 / 1.1\\
 & B4 probe median & 0.1444 (0.0075) & 0.669 & q46: 34.4 / 1.1\\
\addlinespace
ibm\_fez & B1 probe best & 0.2306 (0.0090) & 0.385 & q11: 9.3 / 8.3\\
 & B2 calibration best & 0.2665 (0.0090) & 0.580 & q83: 12.3 / 10.2\\
 & B3 stock level 3 & 0.3413 (0.0124) & 0.592 & q113: 12.1 / 4.0\\
 & B4 probe median & 0.2069 (0.0114) & 0.526 & q83: 12.3 / 10.2\\
\addlinespace
ibm\_kingston & B1 probe best & 0.3392 (0.0124) & 0.546 & q127: 5.0 / 4.5\\
 & B2 calibration best & 0.3417 (0.0098) & 0.638 & q107: 19.8 / 1.0\\
 & B3 stock level 3 & 0.3409 (0.0131) & 0.624 & q107: 19.8 / 1.0\\
 & B4 probe median & 0.2818 (0.0102) & 0.597 & q107: 19.8 / 1.0\\}
\newcommand{\ContrastRows}{ibm\_marrakesh & L1 & B1$-$B2 & +25.02 & [+24.16, +25.89] & $6.8\times10^{-13}$ & $7.8\times10^{-3}$\\
 & L2 & B1$-$B3 & +23.97 & [+23.24, +24.70] & $3.0\times10^{-13}$ & $7.8\times10^{-3}$\\
 & L3 & B1$-$B4 & +16.89 & [+16.10, +17.67] & $6.6\times10^{-12}$ & $7.8\times10^{-3}$\\
 & L4 & (B1$-$B2)$_{\rm mit}$ & -10.15 & [-16.51, -3.79] & $5.7\times10^{-3}$ & $9.8\times10^{-3}$\\
\addlinespace
ibm\_fez & L1 & B1$-$B2 & -3.60 & [-4.37, -2.83] & $4.6\times10^{-6}$ & $7.8\times10^{-3}$\\
 & L2 & B1$-$B3 & -11.07 & [-12.23, -9.90] & $1.4\times10^{-8}$ & $7.8\times10^{-3}$\\
 & L3 & B1$-$B4 & +2.37 & [+1.29, +3.45] & $7.9\times10^{-4}$ & $7.8\times10^{-3}$\\
 & L4 & (B1$-$B2)$_{\rm mit}$ & -19.58 & [-21.26, -17.89] & $3.2\times10^{-9}$ & $7.8\times10^{-3}$\\
\addlinespace
ibm\_kingston & L1 & B1$-$B2 & -0.25 & [-1.62, +1.13] & 1.00 & 1.00\\
 & L2 & B1$-$B3 & -0.17 & [-1.61, +1.27] & 1.00 & 1.00\\
 & L3 & B1$-$B4 & +5.75 & [+4.60, +6.89] & $5.0\times10^{-6}$ & $7.8\times10^{-3}$\\
 & L4 & (B1$-$B2)$_{\rm mit}$ & -9.16 & [-11.68, -6.64] & $5.3\times10^{-5}$ & $7.8\times10^{-3}$\\}
\newcommand{\JobRows}{ibm\_marrakesh & \texttt{dc0apq2nj2vs73ddo74g} & \texttt{dc0aq0qnj2vs73ddo7ig} & 42\\
ibm\_fez & \texttt{dc0aqg5s2omc73891he0} & \texttt{dc0as15s2omc73891lgg} & 39\\
ibm\_kingston & \texttt{dc0av1eq4uuc73a2bk0g} & \texttt{dc0av8uq4uuc73a2bkd0} & 40\\}
\newcommand{\TotalSeconds}{121}
\newcommand{\HelpsCount}{1}


\title{\bfseries A per-job readout probe rescues qubit selection on one IBM Heron r2 processor and
hurts it on another: a preregistered three-device test}
\author{Devin Phillip Davis\\ \small Agile Defense Systems LLC}
\date{20 November 2026 \quad (TESSERACT-MAP-04 results, version 1.0)}

\begin{document}
\maketitle

\begin{abstract}
\noindent
In TESSERACT-MAP-03, a single qubit with a readout fault absent from IBM's published calibration
reversed a 10\,pp layout effect. We preregistered a direct test of the obvious remedy: measure readout
on every qubit in a few-second probe job just before the experiment, and rank 40-qubit paths with the
probe's readout errors instead of the published ones. On \texttt{ibm\_marrakesh}, \texttt{ibm\_fez}
and \texttt{ibm\_kingston} (six jobs, \TotalSeconds{}\,s QPU in total), we compared four layouts of the
same 39-CZ circuit: the probe-ranked best path, the calibration-ranked best path, Qiskit's stock
level-3 layout, and the probe-ranked median path. On \texttt{ibm\_marrakesh} the probe path won by
$+25.0$\,pp (classical fidelity 0.313 against 0.063): the calibration-ranked path again contained
qubit~46, which the probe measured at 34\% readout error while IBM listed 1.1\%. On \texttt{ibm\_fez}
the probe path lost by 3.6\,pp, and the stock layout beat every other arm. On \texttt{ibm\_kingston}
there was no difference, although the calibration path contained a qubit read out at 20\% (listed at
1.0\%). After readout mitigation, the calibration-ranked path scored higher on all three devices. The
probe therefore trades readout quality for gate and coherence quality, and that trade pays only when
a readout fault is severe. By the preregistered criterion (a significant gain on at least two of three
devices), probe-guided selection does not generalise. Hidden readout faults, however, were present on
two of the three processors.
\end{abstract}

\section{Background}
TESSERACT-MAP-02~\cite{map02} found that, at a fixed 39~CZ and depth~66, the choice of 40 physical
qubits changed a circuit's classical fidelity $F_{\mathrm{cl}}$ by 10.5\,pp on \texttt{ibm\_marrakesh}.
Four days later, MAP-03~\cite{map03} found the same layouts reversed: qubit~46 on the
calibration-ranked best path reported~1 for $|0\rangle$ in 54\% of shots, against a listed error near
1\%. If readout faults like this are common and undisclosed, a cheap readout probe before each job
should choose better qubits. MAP-04 tests that idea, preregistered~\cite{prereg} before any job ran.

\section{Design}
Each backend received two jobs, minutes apart.

\paragraph{Probe.} All-$|0\rangle$ and all-$|1\rangle$ on every qubit, 3 repeats each, 2{,}048 shots
(6~PUBs, about 5\,s). A qubit's probe readout error is the mean of $p(1|0)$ and $p(0|1)$.

\paragraph{Selection.} This is MAP-02's path rule: a greedy search over 400 simple 40-qubit paths, with
cost $-\sum\log(1-e_{\mathrm{CZ}})-\sum\log(1-e_{\mathrm{RO}})$ and CZ errors from IBM's calibration
at probe time. The probe paths use the probe's $e_{\mathrm{RO}}$; the calibration paths use IBM's.

\paragraph{Main job.} The 39-CZ chain circuit of MAP-02 (exact four-string output, optimisation
level~3) on four layouts: B1, the probe-ranked best path; B2, the calibration-ranked best path; B3,
Qiskit's stock level-3 layout; B4, the probe-ranked median path. Union all-0/all-1 circuits cover
every qubit used. There were 10 shuffled blocks $\times$ 6 PUBs $\times$ 2{,}048 shots. Before
submission, every arm was checked in a noiseless matrix-product-state simulation
($F_{\mathrm{cl}}\ge0.998$).

\paragraph{Contrasts} (per device; paired $t$ over 10 blocks; Holm across four): L1 = B1$-$B2,
L2 = B1$-$B3, L3 = B1$-$B4, and L4 = B1$-$B2 on readout-mitigated $F_{\mathrm{cl}}$ (tensored inverse,
as in MAP-03). Decision: probe selection \emph{helps} on a device if L1~$>0$ with Holm $p<0.05$, and
\emph{generalises} if it helps on at least two of three devices.

\section{Results}

\begin{table}[t]
\centering\scriptsize\setlength{\tabcolsep}{2pt}
\caption{Jobs and QPU seconds charged (probe + main).}
\begin{tabular}{@{}lllr@{}}
\toprule
Backend & Probe job & Main job & s\\
\midrule
\JobRows
\bottomrule
\end{tabular}
\end{table}

\begin{table*}[t]
\centering\footnotesize
\caption{Arms. $F_{\mathrm{cl}}$: mean over 10 blocks (SD). ``Mitigated'': readout-mitigated
$F_{\mathrm{cl}}$. Last column: the layout's worst qubit by probe readout error, with the probe's and
IBM's readout error (\%).}
\label{tab:arms}
\begin{tabular}{@{}lllrl@{}}
\toprule
Backend & Arm & $F_{\mathrm{cl}}$ (SD) & Mitigated & Worst qubit: probe / IBM (\%)\\
\midrule
\ArmRows
\bottomrule
\end{tabular}
\end{table*}

\begin{table*}[t]
\centering\footnotesize
\caption{Preregistered contrasts (pp of $F_{\mathrm{cl}}$), Holm across four per device; exact
Wilcoxon (Holm) as sensitivity.}
\label{tab:contrasts}
\begin{tabular}{@{}lllrcrr@{}}
\toprule
Backend & ID & Contrast & Mean (pp) & 95\% CI (pp) & Holm $p$ & Wilcoxon $p$\\
\midrule
\ContrastRows
\bottomrule
\end{tabular}
\end{table*}

\paragraph{ibm\_marrakesh.} q46 was still faulty, four days after MAP-03: probe readout error 34\%,
listed at 1.1\%. The calibration rule placed the circuit on it, as did the stock layout and the probe
median path. The probe path avoided it and scored 0.313, against 0.063 for the calibration path
(L1~$=+25.0$\,pp) and 0.074 for the stock layout (L2~$=+24.0$\,pp).

\paragraph{ibm\_fez.} Readout faults there were milder and mostly visible in IBM's own figures. The
probe path lost to the calibration path (L1~$=-3.6$\,pp) and to the stock layout by 11.1\,pp. The
stock layout was the best arm on this device.

\paragraph{ibm\_kingston.} The calibration path contained q107, read out at 19.8\% (listed at 1.0\%),
yet B1 and B2 did not differ (L1~$=-0.25$\,pp, 95\% CI $[-1.62,+1.13]$). What B2 lost on readout it
made up elsewhere.

\paragraph{Mitigated endpoint.} L4 was negative on all three devices ($-10.2$, $-19.6$ and $-9.2$\,pp).
Once readout is corrected, the calibration-ranked path is the better path everywhere. Avoiding a bad
readout qubit forces the 40-qubit path onto couplers and qubits that are worse in other respects.

\paragraph{Path quality.} L3 (probe best vs probe median) was positive on all three devices
($+16.9$, $+2.4$ and $+5.8$\,pp), so the size of the path effect itself is device dependent.

\section{Discussion}
The preregistered verdict is that probe-guided selection helped on \HelpsCount{} of 3 devices and
does not generalise. Two findings stand out nonetheless.

\textbf{Undisclosed readout faults are common.} The two processors with large faults had a qubit
at 20--34\% readout error that IBM's calibration listed at about 1\%. On \texttt{ibm\_marrakesh} the
same qubit has been faulty for at least four days, and every calibration-driven layout we tried, the
stock Qiskit pass included, used it.

\textbf{Correcting the readout term is not enough.} The path cost weighs a qubit's readout error the
same as one CZ error. A 40-qubit path, however, uses 39 CZ gates and many idle periods, and the
probe-ranked paths were worse on those. The remedy that the mitigated endpoint points to is to keep the
calibration-ranked path when its readout faults can be corrected, or to exclude only grossly faulty
qubits. Testing that is a natural next step.

\paragraph{Limitations.} One run per device, on one day. The probe measures only all-0/all-1 readout.
The selection rule is a greedy search over 400 candidates. Mitigation is tensored. B1 and B2 differ in
many qubits at once, so we do not attribute the differences to individual qubits.

\section*{Data and code availability}
The record contains, for each device, the probe counts, the probe-time calibration, the lock, the
QPY, the raw counts, the analysis, and the code (\texttt{phase3/tesseract\_map\_04.py},
\texttt{generate\_tables.py}) that produces every number here, with a SHA-256 manifest.

\section*{Acknowledgements}
We acknowledge the use of IBM Quantum services; the views expressed are the author's and do not
reflect the official policy or position of IBM or the IBM Quantum team. Analysis code and this
manuscript were prepared with the assistance of AI coding tools under the author's direction and
review.

\begin{thebibliography}{9}\small
\bibitem{map02} D.~P.~Davis, Zenodo, doi:10.5281/zenodo.24406426 (2026).
\bibitem{map03} D.~P.~Davis, Zenodo, doi:10.5281/zenodo.24406525 (2026).
\bibitem{prereg} D.~P.~Davis, ``Preregistration: TESSERACT-MAP-04,'' Zenodo, doi:10.5281/zenodo.24406640 (2026).
\end{thebibliography}
\end{document}
